\renewcommand{\figdir}{figures} % figure directory

\plot{sources}{width=0.8\textwidth}{
	{2D schematic representation of sources located in subsurface, over the
	domain $\overline{\mathcal{V}} = \bigcup_{\nu = 1}^{N_{\nu}}$, $N_{\nu} = 2$, 
	with top limited by the plane $z^{t}_{0}$. 
	Points within the sources are represented by $\mathbf{r}' = (x', y', z')^{\top}$.
	The dashed line represents the surface containing the observation points 
	$\mathbf{r} = (x, y, z)^{\top}$. The thick line represents the Earth's surface.
	The planar surface $\mathcal{S}^{\dagger}$ containing points 
	$\mathbf{r}^{\dagger} = (x^{\dagger}, y^{\dagger}, z^{\dagger}_{0})^{\top}$ 
	above the sources is represented by a continuous straight line.
	The fine dashed lines represent two possible definitions of a horizontal surface $\mathcal{S}''$ 
	at a constant vertical coordinate $z''_{0}$ satisfying $z''_{0} > z^{t}_{0}$ or $z''_{0} < z^{t}_{0}$. 
	}
	\label{fig:geometry}
}

\plot{model_view_1.png}{width=0.8\textwidth}{
	{Simulated magnetic source formed by horizontally juxtaposed rectangular prisms.}
	\label{fig:model}
}

\plot{discrete_layers.png}{width=0.8\textwidth}{
	{
		Discrete equivalent layers estimated by running Algorithm \ref{alg:numerical-eqlayer}.
		(a) Computed layer having the constant magnetization direction equal to that of the 
		model shown in Figure \ref{fig:model}.
		(b) Computed layer having the constant magnetization direction different from that of the 
		model shown in Figure \ref{fig:model}.
		All values are in $\mu$T m.	
		The black line represents the projection of the model on the $xy$ plane.	
	}
	\label{fig:discrete-layers}
}

\plot{Neumann_bx_by_bz.png}{width=0.8\textwidth}{
	{
		Potential fields produced by the model shown in Figure \ref{fig:model} and by the discrete 
		layer shown in Figure \ref{fig:discrete-layers}a.
		(a), (d) and (g) show the potential fields produced by the discrete layer.
		(b), (e) and (h) show the potential fields produced by the model.
		(c), (f) and (i) show the difference between the fields produced by the layer and the model.
		The upper, middle and bottom rows show the $x$-, $y$- and $z$-components of the 
		magnetic induction field, respectively. All values are in nT.
		The black line represents the projection of the model on the $xy$ plane.	
	}
	\label{fig:Neumann-bx-by-bz}
}

\plot{Neumann_atfa_b_rtp.png}{width=0.8\textwidth}{
	{
		Potential fields produced by the model shown in Figure \ref{fig:model} and by the discrete 
		layer shown in Figure \ref{fig:discrete-layers}a.
		(a), (d) and (g) show the potential fields produced by the discrete layer.
		(b), (e) and (h) show the potential fields produced by the model.
		(c), (f) and (i) show the difference between the fields produced by the layer and the model.
		The upper, middle and bottom rows show the approximated total-field anomaly, amplitude of 
		magnetic induction field and RTP anomaly, respectively. All values are in nT.
		The black line represents the projection of the model on the $xy$ plane.		
	}
	\label{fig:Neumann-atfa-b-rtp}
}

\plot{generic_bx_by_bz.png}{width=0.8\textwidth}{
	{
		Potential fields produced by the model shown in Figure \ref{fig:model} and by the discrete 
		layer shown in Figure \ref{fig:discrete-layers}b.
		(a), (d) and (g) show the potential fields produced by the discrete layer.
		(b), (e) and (h) show the potential fields produced by the model.
		(c), (f) and (i) show the difference between the fields produced by the layer and the model.
		The upper, middle and bottom rows show the $x$-, $y$- and $z$-components of the 
		magnetic induction field, respectively. All values are in nT.
		The black line represents the projection of the model on the $xy$ plane.	
	}
	\label{fig:dipole-bx-by-bz}
}

\plot{generic_atfa_b_rtp.png}{width=0.8\textwidth}{
	{
		Potential fields produced by the model shown in Figure \ref{fig:model} and by the discrete 
		layer shown in Figure \ref{fig:discrete-layers}b.
		(a), (d) and (g) show the potential fields produced by the discrete layer.
		(b), (e) and (h) show the potential fields produced by the model.
		(c), (f) and (i) show the difference between the fields produced by the layer and the model.
		The upper, middle and bottom rows show the approximated total-field anomaly, amplitude of 
		magnetic induction field and RTP anomaly, respectively. All values are in nT.
		The black line represents the projection of the model on the $xy$ plane.		
	}
	\label{fig:dipole-atfa-b-rtp}
}

\plot{generic_rtp_fixed.png}{width=0.8\textwidth}{
	{
		RTP anomalies produced by the model shown in Figure \ref{fig:model} and by the discrete 
		layer shown in Figure \ref{fig:discrete-layers}b, after correction.
		(a) shows the fixed RTP anomaly produced by the discrete layer.
		(b) shows the RTP anomaly produced by the model.
		(c) shows the difference between the RTP anomalies produced by the layer and the model.
		All values are in nT.
		The black line represents the projection of the model on the $xy$ plane.		
	}
	\label{fig:dipole-rtp-fixed}
}

